\documentclass[11pt]{article}
\usepackage[T1]{fontenc}
\usepackage{textcomp}
\usepackage[margin=0.75in]{geometry}
\usepackage{amsmath,amssymb,amsthm}
\usepackage{array,booktabs}
\usepackage{hyperref}
\setlength{\parindent}{0pt}
\newtheorem{theorem}{Theorem}
\theoremstyle{definition}
\newtheorem{definition}{Definition}
\title{Flow Proof Artifact: prop-04-square-on-whole-cut-random}
\author{Generated by \texttt{flow doc proof}}
\date{Flow Proof Book}
\begin{document}
\maketitle
If a straight line is cut at random, the square on the whole equals the squares on the segments together with twice the rectangle contained by the segments.\\[0.5em]
\textbf{Source.} euclid — Elements, Book II, Proposition 4\\[0.5em]
\bigskip
\noindent{\large \textbf{Derived fact 1.} \textit{If a straight line is cut at random, the square on the whole equals the squares on the segments together with twice the rectangle contained by the segments} \hfill the Euclidean plane \cdot Euclid Book II \cdot Proposition 4: the square on the whole equals the squares on the parts plus twice their rectangle}\par
\smallskip\noindent\textit{Coordinate: the Euclidean plane \cdot Euclid Book II \cdot Proposition 4: the square on the whole equals the squares on the parts plus twice their rectangle}\par
\smallskip\noindent\textit{Source: euclid — Elements, Book II, Proposition 4}\par
\smallskip\noindent\textit{Built on: proposition 2: the rectangle by the whole and one part equals the parts plus a square, for Euclid Book II on the Euclidean plane, proposition 3: the rectangle by the whole and the other part equals the parts plus a square, for Euclid Book II on the Euclidean plane}\par
\medskip
\noindent\textbf{Goal.} If a straight line is cut at random, the square on the whole equals the squares on the segments together with twice the rectangle contained by the segments\par
\noindent$\displaystyle sq(whole) = sq(C) + sq(D) + 2 \cdot rect(C, D)$\par
\medskip
\noindent
\begin{tabular}{@{} >{\raggedright\arraybackslash}p{0.06\textwidth} >{\raggedright\arraybackslash}p{0.47\textwidth} >{\raggedleft\arraybackslash}p{0.41\textwidth} @{}}
\textbf{\#} & \textbf{Proof} & \textbf{Mathematics} \\
\midrule
\textbf{1.} & We prove that proposition 4: the square on the whole equals the squares on the parts plus twice their rectangle for Euclid Book II the Euclidean plane. & \\[0.45em]
\textbf{2.} & Let whole = straight\_line\_AB. & \\[0.45em]
\textbf{3.} & Let C = segment\_AC. & \\[0.45em]
\textbf{4.} & Let D = segment\_CB. & \\[0.45em]
\textbf{5.} & From step 2, step 3, and step 4, this implies sq(whole) equals sq(C) plus sq(D) plus 2 times rect(C, D). Hence proven. & $\displaystyle sq(whole) = sq(C) + sq(D) + 2 \cdot rect(C, D)$ \\[0.45em]
\textbf{6.} & We invoke the derived fact governing Euclid Book II the Euclidean plane: proposition 2: the rectangle by the whole and one part equals the parts plus a square, for Euclid Book II on the Euclidean plane. & \\[0.45em]
\textbf{7.} & We invoke the derived fact governing Euclid Book II the Euclidean plane: proposition 3: the rectangle by the whole and the other part equals the parts plus a square, for Euclid Book II on the Euclidean plane. & \\[0.45em]
\bottomrule
\end{tabular}
\label{thm:Geometry-Euclid Book II-proposition_4_the_square_on_the_whole_equals_the_squares_on_the_parts_plus_twice_their_rectangle}
\par\medskip\hrule
\end{document}
